% Generated draft: requires mathematical and source review.
\documentclass[11pt]{article}
\usepackage[margin=25mm]{geometry}
\usepackage{fontspec}
\setmainfont{Latin Modern Roman}
\usepackage{amsmath,amssymb,mathtools,unicode-math}
\setmathfont{Latin Modern Math}
\usepackage{enumitem,needspace,xurl,hyperref}
\setlength{\parindent}{0pt}
\setlength{\parskip}{5pt}
\setlength{\emergencystretch}{3em}
\newcommand{\N}{\mathbb{N}}
\newcommand{\Z}{\mathbb{Z}}
\newcommand{\Q}{\mathbb{Q}}
\newcommand{\R}{\mathbb{R}}
\newcommand{\C}{\mathbb{C}}
\newcommand{\F}{\mathbb{F}}
\newcommand{\A}{\mathbb{A}}
\newcommand{\PP}{\mathbb{P}}
\newcommand{\E}{\mathbb{E}}
\newcommand{\eps}{\varepsilon}
\newcommand{\dd}{\,\mathrm{d}}
\begin{document}
\section*{UnsolvedMath ID: 30004395}
\textbf{Source problem code:} catalogue entry
\subsection{Definitions and mathematical statement}
Let \(M\) be an oriented, closed, connected manifold of dimension \(n\), where closed means compact without boundary. No dimension restriction is supplied. The manifold is aspherical if its universal covering space is contractible. For \(A\in\{\Z,\R\}\), write \([M]_A\in H_n(M;A)\) for its fundamental class. For a finite singular chain written with distinct singular simplices as \(c=\sum_{j=1}^r a_j\sigma_j\), define \(\lVert c\rVert_1=\sum_{j=1}^r|a_j|\). Define the real and integral simplicial volumes by
\[
\lVert M\rVert=\inf\{\lVert c\rVert_1:c\in C_n(M;\R),\ \partial c=0,\ [c]=[M]_{\R}\},
\]
\[
\lVert M\rVert_{\Z}=\inf\{\lVert c\rVert_1:c\in C_n(M;\Z),\ \partial c=0,\ [c]=[M]_{\Z}\}.
\]
The stable integral simplicial volume is
\[
\lVert M\rVert_{\Z}^{\infty}
=\inf_{p:N\to M}\frac{\lVert N\rVert_{\Z}}{\deg(p)},
\]
where the infimum ranges over all connected finite-sheeted covering maps, including the identity covering, and \(N\) carries the orientation induced from \(M\), so that \(\deg(p)\) is the positive number of sheets. Write
\[
\chi(M)=\sum_{k=0}^{n}(-1)^k\dim_{\Q}H_k(M;\Q).
\]
A group \(G\) is residually finite if, for every \(g\in G\setminus\{1\}\), there is a homomorphism \(q:G\to F\) to a finite group such that \(q(g)\ne1\).

\textbf{First question.} For every oriented, closed, connected, aspherical manifold \(M\), is the implication
\[
\lVert M\rVert=0\quad\Longrightarrow\quad\chi(M)=0
\]
valid?

\textbf{Integral approximation question.} For every such \(M\) satisfying \(\lVert M\rVert=0\) and having residually finite fundamental group, does
\[
\lVert M\rVert_{\Z}^{\infty}=\lVert M\rVert=0
\]
hold?

The supplied source describes the second question as ``more strongly.'' Its vanishing hypothesis is retained from the opening sentence; no assertion about integral approximation for manifolds of positive simplicial volume is added. The source does not specify a manifold category or an explicit dimension range, and these qualifications remain unresolved.
\subsection{Short English statement}
Must every oriented, closed, connected, aspherical manifold with zero real simplicial volume have zero Euler characteristic? Under the additional assumption that its fundamental group is residually finite, must its stable integral simplicial volume also vanish? Both questions are supplied as open; this draft does not independently verify their status.
\subsection{Sources}
\begin{itemize}
\item Local UnsolvedMath registry, supplied complete record with permanent ID \(30004395\), source problem code \texttt{catalogue entry}, title \emph{Euler Characteristic and Integral Approximation from Simplicial Volume}, in \emph{Oberwolfach Reports Open Problems}. The record gives proposed year 2020 and does not name an individual proposer.
\item Supplied source citation: \emph{Manifolds and Groups (2021)}, Oberwolfach Reports, pp.~417--463, DOI \texttt{10.4171/owr/2020/7}. \url{https://doi.org/10.4171/owr/2020/7}. The supplied background describes these reports as documenting workshops and their open-problem discussions. The relationship between the cited year 2021, the DOI's 2020 component, and the proposed year 2020 is not resolved here.
\item Supplied literature source, labelled ``Loh--Moraschini--Raptis 2022'': \url{https://doi.org/10.1007/s40687-022-00341-3}. No article title or passage supporting a specific claim is supplied.
\item Literature assessment dated exactly \texttt{2026-08-22}: supplied dated curation, not independently verified. It labels the problem open and describes Gromov's implication from vanishing simplicial volume to vanishing Euler characteristic for closed aspherical manifolds as a major open problem, with known examples covering only special classes. The supplied list of sources checked contains the Oberwolfach report and the Loh--Moraschini--Raptis reference above. The background explicitly requests independent verification of time-sensitive status and cited claims.
\item The registry labels the statement \texttt{corrected\_verified} and reports that the official report was checked and that both questions from its source passage were retained. These are supplied source-verification claims; this draft does not claim an independent visit to, or verification of, either reference. Attribution to Gromov is supplied for the first implication, without a separate original citation.
\end{itemize}
\end{document}
